% Work/MAT – přepracovaná 18. hodina. Pouze slovní úlohy se třemi množinami.
% Návaznost na 17_MNZ_SlovniUlohyDveMnoziny a 14_MNZ_VennovyDiagramy.
% Učebnice: ucebnice/01_Zakladni_poznatky/Mnoziny/Množiny.tex;
% Realisticky.cz 010306 a 010307: didaktická inspirace (veřejné indexované ukázky).
% Konkrétní zadání a číselné údaje jsou vlastní; lokální PDF lekcí neověřena.
\documentclass[11pt,a4paper]{article}
% Spolecny styl pro ucitelske pripravy MAT (Work/MAT)
% Verze 1.1, 2026-10-23
\NeedsTeXFormat{LaTeX2e}
\ProvidesPackage{priprava}[2026/10/23 v1.1 Sablona ucitelskych priprav MAT]

% Zakladni typografie a cestina
\RequirePackage[T1]{fontenc}
\RequirePackage[utf8]{inputenc}
\RequirePackage[czech]{babel}
\RequirePackage{lmodern}
\RequirePackage{microtype}
\RequirePackage{amsmath,amssymb}
\RequirePackage{xcolor}
\RequirePackage{enumitem}
\RequirePackage[most]{tcolorbox}
\RequirePackage{tabularx,booktabs,array}
\RequirePackage{fancyhdr}
\RequirePackage{needspace}
\RequirePackage{geometry}
\RequirePackage[hidelinks]{hyperref}

\geometry{a4paper,top=17mm,bottom=18mm,left=23mm,right=23mm}
\definecolor{notegray}{gray}{0.39}
\setlength{\parindent}{0pt}
\setlength{\parskip}{2.5pt plus 1pt minus 0.5pt}
\setlist[itemize]{topsep=2pt,itemsep=1pt,parsep=0pt,leftmargin=1.45em}
\setlist[enumerate]{topsep=2pt,itemsep=2pt,parsep=0pt,leftmargin=1.65em}
\setcounter{secnumdepth}{0}
\renewcommand{\arraystretch}{1.12}

% Zahlavi: vlevo rocnik | globalni tema, vpravo cislo: nazev hodiny
% Zapati zustava zcela prazdne.
\pagestyle{fancy}
\fancyhf{}
\renewcommand{\headrulewidth}{0.35pt}
\renewcommand{\footrulewidth}{0pt}
\setlength{\headheight}{14pt}
\newcommand{\nastavPripravu}[4]{%
  \fancyhead[L]{\small #1\enspace\textbar\enspace #2}%
  \fancyhead[R]{\small #3: #4}%
}

% Minimalisticke cernobile boxy. Kazda definice zvlast.
\tcbset{commonlesson/.style={enhanced,breakable=false,colback=white,colframe=black,
  colbacktitle=white,coltitle=black,boxrule=0.45pt,
  arc=0.8mm,outer arc=0.8mm,left=3mm,right=3mm,top=1.6mm,bottom=1.6mm,
  toptitle=1mm,bottomtitle=0.7mm,fonttitle=\bfseries,
  attach boxed title to top left={xshift=2.5mm,yshift=-1.15mm},
  boxed title style={colback=white,colframe=white,boxrule=0pt,left=1mm,right=1mm}}}
\newtcolorbox{zadaniBox}[1][]{commonlesson,title={Zadání},#1}
\newtcolorbox{prikladBox}[1][]{commonlesson,title={Příklad},#1}
\newtcolorbox{definiceBox}[1][]{commonlesson,title={Definice},#1}

% Spolecny box pro zadani a jeho reseni, bez mezery mezi castmi.
% Pouziti:
%   \begin{prikladReseniBox}
%     Zadani ulohy.
%     \tcblower
%     \begin{align*} ... \end{align*}
%   \end{prikladReseniBox}
% Spodni cast ma jemny sedy podklad a delici linku; lze v ni sazet
% vice radku matematiky. Obe casti maji jen jeden spolecny ramecek.
\tcbset{solutionlesson/.style={
  commonlesson,
  enhanced jigsaw,
  breakable,
  colbacklower=black!3,
  segmentation style={solid,draw=black!25,line width=0.35pt},
  fontlower=\small\color{notegray},
  before lower={\noindent\textbf{Řešení:}\enspace},
  leftlower=3mm,rightlower=3mm
}}
\newtcolorbox{prikladReseniBox}[1][]{solutionlesson,title={Příklad},#1}
\newtcolorbox{zadaniReseniBox}[1][]{solutionlesson,title={Zadání},#1}

% Pomucky pro autora pripravy: sedive poznamky; \reseni pro starsi pripravy
\newcommand{\poznamka}[1]{\par\begingroup\small\color{notegray}\itshape
  \textnormal{\textbf{Poznámka pro učitele.}}\enspace #1\par\endgroup}
\newcommand{\reseni}[1]{\par\begingroup\small\color{notegray}
  \textbf{Očekávané řešení:}\enspace #1\par\endgroup}
\newcommand{\blok}[1]{\needspace{5\baselineskip}\subsection*{#1}}
\newcommand{\vyrok}[1]{\uv{#1}}

\usepackage{tikz}
\nastavPripravu{1. ročník}{Množiny}{18}{Slovní úlohy: tři množiny}
\begin{document}

\blok{Od čtyř oblastí k osmi}
V minulé hodině jsme rozlišovali dvě vlastnosti. Se třetí vlastností
má každý prvek jednu z \textbf{osmi možných kombinací} příslušnosti.
Na rozdíl od výčtu prvků dnes do oblastí píšeme \emph{počty}.
Nezačínej u součtů: \textbf{první doplň střed, potom právě dvě oblasti,
právě jednu a nakonec žádnou.}

\begin{definiceBox}
Pro konečné množiny $A,B,C$ znamená $|A\cap B\cap C|$ počet prvků
patřících \textbf{do všech tří}. Naproti tomu $|A\cap B|$ zahrnuje také
prvky, které současně patří do $C$.
\end{definiceBox}

\begin{prikladReseniBox}[title={P1: tři zájmové oblasti – úplný diagram}]
Ve skupině $40$ lidí se zajímá o robotiku $23$, programování $21$
a elektroniku $17$. O robotiku a programování se zajímá $10$,
o robotiku a elektroniku $9$, o programování a elektroniku $8$.
Všechny tři zájmy má $4$ lidí. Urči počty v \emph{osmi oblastech},
počet s \emph{právě dvěma}, \emph{alespoň dvěma} a bez zájmu.
\tcblower
Označme uvedené zájmy $A,B,C$. Nejprve střed $4$.
Právě dvě: $AB$ bez $C$: $10-4=6$;
$AC$ bez $B$: $9-4=5$; $BC$ bez $A$: $8-4=4$.
\[
\begin{aligned}
\text{jen }A&=23-6-5-4=8,\\
\text{jen }B&=21-6-4-4=7,\\
\text{jen }C&=17-5-4-4=4.
\end{aligned}
\]
\begin{center}
\begin{tikzpicture}[scale=.82]
\draw[thick] (-3.06,-2.62) rectangle (3.1,2.04);
\draw[thick] (-.85,.38) circle(1.43);
\draw[thick] (.85,.38) circle(1.43);
\draw[thick] (0,-.83) circle(1.43);
\node at (-1.8,1.72) {$A$};\node at (1.8,1.72) {$B$};
\node at (0,-2.4) {$C$};\node at (2.83,1.78) {$U$};
\node at (-1.45,.59) {$8$};\node at (1.45,.59) {$7$};
\node at (0,1.12) {$6$};\node at (-.82,-.79) {$5$};
\node at (.82,-.79) {$4$};\node at (0,-.23) {$4$};
\node at (0,-1.63) {$4$};\node at (2.55,-2.09) {$2$};
\end{tikzpicture}
\end{center}
Právě dva zájmy: $6+5+4=15$; alespoň dva: $15+4=19$.
Alespoň jeden: $8+7+4+6+5+4+4=38$; žádný $40-38=2$.
\textbf{Kontrola:} $8+7+4+6+5+4+4+2=40$,
$8+6+5+4=23$ pro skupinu $A$, analogicky pro $B,C$.
\end{prikladReseniBox}

\poznamka{P1 veď společně a na tabuli doplňuj oblasti po etapách.
Žáci při každé etapě počítají sami, ne pouze opisují. Na P1 vyhraď
asi 16–18 minut; kresbu lze předkreslit před začátkem hodiny.}

\blok{Proč musíme společný střed zase přičíst?}
\begin{definiceBox}
\textbf{Inkluze a exkluze pro tři konečné množiny:}
\[
\begin{aligned}
|A\cup B\cup C|={}&|A|+|B|+|C|\\
&-|A\cap B|-|A\cap C|-|B\cap C|+|A\cap B\cap C|.
\end{aligned}
\]
\end{definiceBox}
Součtem $|A|+|B|+|C|$ je střed započten třikrát; odečtením
párových průniků ho zase třikrát odebereme, proto ho jednou vracíme.
V P1 tak $23+21+17-10-9-8+4=38$.
\poznamka{Vzorec je zkratka a kontrola diagramu, nikoliv náhrada jeho vyplnění.
Pro první seznámení stačí slovní odůvodnění a porovnání s P1.}

\blok{Samostatné procvičení: další situace}
\begin{zadaniReseniBox}[title={S1: tři druhy vad výrobků}]
V sérii $45$ výrobků bylo zjištěno $24$ s vadou $A$, $22$ s vadou $B$
a $19$ s vadou $C$. Vady $A$ i $B$ má $10$ výrobků,
$A$ i $C$ má $9$, $B$ i $C$ má $8$; všechny tři vady mají $4$ výrobky.
Urči všech osm oblastí, počet výrobků s právě dvěma vadami a počet
bez těchto tří vad. Ověř výsledek druhým způsobem.
\tcblower
Střed $4$; právě dva typy vad: $AB$ pouze $6$, $AC$ pouze $5$,
$BC$ pouze $4$. Pouze $A$: $24-6-5-4=9$;
pouze $B$: $22-6-4-4=8$; pouze $C$: $19-5-4-4=6$.
Sjednocení $9+8+6+6+5+4+4=42$; bez těchto vad $45-42=3$.
Právě dvě vady $6+5+4=15$.
Kontrola vzorcem: $24+22+19-10-9-8+4=42$.
\end{zadaniReseniBox}

\begin{zadaniReseniBox}[title={S2: od celku k jednotlivým oblastem},breakable=false]
V $48$členné skupině platí $|A|=24$, $|B|=21$, $|C|=18$,
$|A\cap B|=10$, $|A\cap C|=8$, $|B\cap C|=7$
a $|A\cap B\cap C|=3$. Kolik lidí není v žádné skupině?
Kolik je právě v jedné a právě ve dvou?
\tcblower
Jen $AB$ $7$, jen $AC$ $5$, jen $BC$ $4$; pouze $A$ $9$,
pouze $B$ $7$, pouze $C$ $6$, střed $3$.
Právě v jedné: $9+7+6=22$; právě ve dvou $7+5+4=16$.
Ve sjednocení $22+16+3=41$, žádná $48-41=7$.
\end{zadaniReseniBox}
\poznamka{Jádro: 4 min návrat ke čtyřem oblastem, 17 min P1,
6 min odvození vzorce, 13 min samostatná S1, 5 min kontrola a závěr.
S2 je rezerva nebo domácí dokončení. Celkem 45 minut.}

\blok{Volitelná rezerva: slovní významy}
\begin{zadaniReseniBox}[title={R1: nejvýše dvě skupiny},breakable=false]
Pro S1 vypočti, kolik výrobků má nejvýše dvě z uvedených vad
a kolik alespoň jednu. Mohou dvě párové oblasti být stejné?
\tcblower
Nejvýše dvě vady: $45-4=41$; alespoň jednu $42$.
Stejné počty v různých oblastech být mohou, i když jde o různé výrobky.
\end{zadaniReseniBox}

\blok{Domácí procvičování}
\begin{zadaniReseniBox}[title={D1: diagram vždy od středu}]
Ve skupině $36$ osob má vlastnost $A$ $20$ lidí, $B$ $18$, $C$ $16$;
$|A\cap B|=9$, $|A\cap C|=8$, $|B\cap C|=7$ a
$|A\cap B\cap C|=4$.
\begin{enumerate}[label=\alph*)]
\item Vyplň osm oblastí diagramu.
\item Kolik lidí má právě dvě vlastnosti a kolik alespoň dvě?
\item Kolik lidí nemá žádnou z vlastností? Ověř vzorcem.
\end{enumerate}
\tcblower
Střed $4$, $AB$ bez $C$ $5$, $AC$ bez $B$ $4$, $BC$ bez $A$ $3$;
pouze $A$ $7$, pouze $B$ $6$, pouze $C$ $5$, mimo všechny $2$.
Právě dvě: $12$, alespoň dvě: $16$.
Ve sjednocení $20+18+16-9-8-7+4=34$, mimo všechny $36-34=2$.
\end{zadaniReseniBox}
\end{document}
